\documentclass[12pt,a4paper]{article}
\usepackage[T1]{fontenc}
\usepackage{amsmath,amssymb,mathrsfs,tikz,times,pifont}
\usepackage{enumitem}
\usepackage{multicol}
\usepackage{lmodern}
\usetikzlibrary{trees}
\newcommand\circitem[1]{%
\tikz[baseline=(char.base)]{
\node[circle,draw=gray, fill=red!55,
minimum size=1.2em,inner sep=0] (char) {#1};}}
\newcommand\boxitem[1]{%
\tikz[baseline=(char.base)]{
\node[fill=cyan,
minimum size=1.2em,inner sep=0] (char) {#1};}}
\setlist[enumerate,1]{label=\protect\circitem{\arabic*}}
\setlist[enumerate,2]{label=\protect\boxitem{\alph*}}
%%%::::::by chnini ameur :::::::%%%
\everymath{\displaystyle}
\usepackage[left=1cm,right=1cm,top=1cm,bottom=1.7cm]{geometry}
\usepackage[colorlinks=true, linkcolor=red, urlcolor=red, citecolor=red]{hyperref}
\usepackage{array,multirow}
\usepackage[most]{tcolorbox}
\usepackage{varwidth}
\usepackage{float} %pour utiliser l'option [H] qui force l'image à apparaître exactement à l'endroit où elle est placée dans le code.
\tcbuselibrary{skins,hooks}
\usetikzlibrary{patterns}
%%%::::::by chnini ameur :::::::%%%
\newtcolorbox{exa}[2][]{enhanced,breakable,before skip=2mm,after skip=5mm,
colback=yellow!20!white,colframe=black!20!blue,boxrule=0.5mm,
attach boxed title to top left ={xshift=0.6cm,yshift*=1mm-\tcboxedtitleheight},
fonttitle=\bfseries,
title={#2},#1,
% varwidth boxed title*=-3cm,
boxed title style={frame code={
\path[fill=tcbcolback!30!black]
([yshift=-1mm,xshift=-1mm]frame.north west)
arc[start angle=0,end angle=180,radius=1mm]
([yshift=-1mm,xshift=1mm]frame.north east)
arc[start angle=180,end angle=0,radius=1mm];
\path[left color=tcbcolback!60!black,right color = tcbcolback!60!black,
middle color = tcbcolback!80!black]
([xshift=-2mm]frame.north west) -- ([xshift=2mm]frame.north east)
[rounded corners=1mm]-- ([xshift=1mm,yshift=-1mm]frame.north east)
-- (frame.south east) -- (frame.south west)
-- ([xshift=-1mm,yshift=-1mm]frame.north west)
[sharp corners]-- cycle;
},interior engine=empty,
},interior style={top color=yellow!5}}
%%%%%%%%%%%%%%%%%%%%%%%
\usepackage{fancyhdr}
\usepackage{eso-pic}         % Pour ajouter des éléments en arrière-plan
% Commande pour ajouter du texte en arrière-plan
\usepackage{tkz-tab}
% Filigrane "PGB" centré sur chaque page
\AddToShipoutPicture{
  \AtTextCenter{%
    \makebox[0pt]{\rotatebox{50}{\textcolor[gray]{0.6}{\fontsize{1.5cm}{1.5cm}\selectfont maths-mes-maths}}}
  }
}
\usepackage{lastpage}
\fancyhf{}
\pagestyle{fancy}
\renewcommand{\footrulewidth}{1pt}
\renewcommand{\headrulewidth}{0pt}
\renewcommand{\footruleskip}{10pt}
\fancyfoot[R]{
\color{red}\ding{45}\ \textbf{2026}
}
\fancyfoot[L]{
\color{red}\ding{45}\ \textbf{Prof:M. Diatta}
}
\cfoot{\bf
\thepage /
\pageref{LastPage}}
% Création du compteur pour les exercices
\newcounter{exercice}
\renewcommand{\theexercice}{\arabic{exercice}}  % Définit l'affichage du compteur en chiffres arabes

% Définir la commande \exo
\newcommand{\exo}{\refstepcounter{exercice}\textbf{Exercice \theexercice} }
\begin{document}
\renewcommand{\arraystretch}{1.5}
\renewcommand{\arrayrulewidth}{1.2pt}
\begin{tcolorbox}[
    colback=white,
    colframe=blue!60!black,
    boxrule=0.8pt,
    arc=3pt,
    left=3mm, right=3mm, top=2.5mm, bottom=2.5mm
]
    \begin{center}
        \textbf{\large Lycée de Diattacounda} \hfill \textbf{A.S. : 2026/2027} \\[1mm]
        \textbf{Mathématiques} $\bullet$ \textbf{Niveau : TS2} \hfill \textbf{Date : 13/11/2026} \\[2mm]
        \rule{0.8\linewidth}{0.4pt} \\[1.5mm]
        {\Large\bfseries\sffamily TD : Similitudes}
    \end{center}
\end{tcolorbox}
\vspace{4mm}


\fbox{\textbf{\exo}} 
Soit $s$ la similitude directe de centre $\Omega$ d'affixe $\omega$, de rapport $k$ et d'angle $\theta$.
Son écriture complexe est

\[
z' - \omega = k\,e^{i\theta}\,(z-\omega)
\iff z' = az + b,
\quad \text{avec} \quad a = k\,e^{i\theta} \ \text{ et } \ b = (1-a)\,\omega .
\]

\begin{multicols}{2}
\setlength{\columnseprule}{0.1mm} % La largeur de la ligne verticale entre les colonnes

\begin{enumerate}

    \item $\omega = 0$, $k=2$, $\theta=\dfrac{\pi}{6}$.

    $
    \begin{aligned}
    a &= 2\,e^{i\pi/6} \\[1mm]
      &= 2\left(\dfrac{\sqrt3}{2}+\dfrac{i}{2}\right) \\[1mm]
      &= \sqrt3 + i,
    \end{aligned}
    $

    $
    \begin{aligned}
    b &= (1-a)\,\omega \\[1mm]
      &= (1-a)\times 0 \\[1mm]
      &= 0.
    \end{aligned}
    $

    \[
    \boxed{z' = (\sqrt3 + i)\,z}
    \]

    \item $\omega = 1$, $k=\sqrt2$, $\theta=\dfrac{\pi}{2}$.

    $
    \begin{aligned}
    a &= \sqrt2\,e^{i\pi/2} \\[1mm]
      &= \sqrt2\,(0 + i) \\[1mm]
      &= i\sqrt2,
    \end{aligned}
    $

    $
    \begin{aligned}
    b &= (1-a)\,\omega \\[1mm]
      &= \left(1 - i\sqrt2\right)\times 1 \\[1mm]
      &= 1 - i\sqrt2.
    \end{aligned}
    $

    \[
    \boxed{z' = i\sqrt2\,z + 1 - i\sqrt2}
    \]

    \item $\omega = 2 - i$, $k=1$, $\theta=\dfrac{5\pi}{6}$.

    $
    \begin{aligned}
    a &= e^{5i\pi/6} \\[1mm]
      &= \cos\dfrac{5\pi}{6} + i\sin\dfrac{5\pi}{6} \\[1mm]
      &= -\dfrac{\sqrt3}{2}+\dfrac{i}{2},
    \end{aligned}
    $

    $
    \begin{aligned}
    b &= (1-a)\,\omega \\[1mm]
      &= \left(1+\dfrac{\sqrt3}{2}-\dfrac{i}{2}\right)(2-i) \\[1mm]
      &= 2 + \sqrt3 - i\left(1+\dfrac{\sqrt3}{2}\right) - i - \dfrac12 \\[1mm]
      &= \dfrac32+\sqrt3 - i\left(2+\dfrac{\sqrt3}{2}\right).
    \end{aligned}
    $

    \[
    \boxed{z' = \left(-\dfrac{\sqrt3}{2}+\dfrac{i}{2}\right)z + \dfrac32+\sqrt3 - i\left(2+\dfrac{\sqrt3}{2}\right)}
    \]

    \item $\omega = -2 + \dfrac{i}{2}$, $k=\dfrac12$, $\theta=-\dfrac{\pi}{3}$.

    $
    \begin{aligned}
    a &= \dfrac12\,e^{-i\pi/3} \\[1mm]
      &= \dfrac12\left(\dfrac12 - \dfrac{i\sqrt3}{2}\right) \\[1mm]
      &= \dfrac14 - \dfrac{i\sqrt3}{4},
    \end{aligned}
    $

    $
    \begin{aligned}
    b &= (1-a)\,\omega \\[1mm]
      &= \left(\dfrac34+\dfrac{i\sqrt3}{4}\right)\left(-2+\dfrac{i}{2}\right) \\[1mm]
      &= -\dfrac32 + \dfrac{3i}{8} - \dfrac{i\sqrt3}{2} - \dfrac{\sqrt3}{8} \\[1mm]
      &= -\dfrac32-\dfrac{\sqrt3}{8} + i\left(\dfrac38-\dfrac{\sqrt3}{2}\right).
    \end{aligned}
    $

    \[
    \boxed{z' = \left(\dfrac14-\dfrac{i\sqrt3}{4}\right)z - \dfrac32-\dfrac{\sqrt3}{8} + i\left(\dfrac38-\dfrac{\sqrt3}{2}\right)}
    \]

\end{enumerate}
\end{multicols}

\fbox{\textbf{\exo}} 

On écrit $f$ sous la forme $z' = az + b$ et on étudie $a$ :
\begin{itemize}
    \item si $a = 1$, $f$ est la \textbf{translation} de vecteur $\vec{u}$ d'affixe $b$ ;
    \item si $a \neq 1$, $f$ admet un unique point invariant $\Omega$ d'affixe $\omega = \dfrac{b}{1-a}$, et
    \[
    z' - \omega = a\,(z-\omega).
    \]
    Alors $f$ est la similitude directe de centre $\Omega$, de rapport $|a|$ et d'angle $\arg(a)$.
    En particulier : si $a \in \mathbb{R}\setminus\{1\}$, c'est une \textbf{homothétie} de rapport $a$ ;
    si $|a| = 1$, c'est une \textbf{rotation}.
\end{itemize}

\begin{enumerate}

    \item $z' = (\sqrt3 + i)\,z$.

    $
    \begin{aligned}
    a &= \sqrt3 + i \\[1mm]
      &= 2\left(\dfrac{\sqrt3}{2}+\dfrac{i}{2}\right) \\[1mm]
      &= 2\,e^{i\pi/6},
    \end{aligned}
    $

    $
    \begin{aligned}
    \omega &= \dfrac{b}{1-a} \\[1mm]
           &= \dfrac{0}{1-a} \\[1mm]
           &= 0.
    \end{aligned}
    $

    \[
    \boxed{f \text{ est la similitude directe de centre } O \text{, de rapport } 2 \text{ et d'angle } \dfrac{\pi}{6}.}
    \]

    \item $z' = -2z + i$.

    $
    \begin{aligned}
    a &= -2 \in \mathbb{R}\setminus\{1\},
    \end{aligned}
    $

    $
    \begin{aligned}
    \omega &= \dfrac{b}{1-a} \\[1mm]
           &= \dfrac{i}{1-(-2)} \\[1mm]
           &= \dfrac{i}{3}.
    \end{aligned}
    $

    \[
    \boxed{f \text{ est l'homothétie de centre } \Omega\left(\dfrac{i}{3}\right) \text{ et de rapport } -2.}
    \]

    \item $z' = -z + 2i$.

    $
    \begin{aligned}
    a &= -1,
    \end{aligned}
    $

    $
    \begin{aligned}
    \omega &= \dfrac{b}{1-a} \\[1mm]
           &= \dfrac{2i}{1-(-1)} \\[1mm]
           &= \dfrac{2i}{2} \\[1mm]
           &= i.
    \end{aligned}
    $

    \[
    \boxed{f \text{ est la symétrie centrale (homothétie de rapport } -1\text{) de centre } \Omega(i).}
    \]

    \item $z' = \dfrac{3 + i\sqrt3}{4}\,z + \dfrac{1 - i\sqrt3}{2}$.

    $
    \begin{aligned}
    a &= \dfrac{3 + i\sqrt3}{4} \\[1mm]
      &= \dfrac{\sqrt3}{2}\left(\dfrac{\sqrt3}{2}+\dfrac{i}{2}\right) \\[1mm]
      &= \dfrac{\sqrt3}{2}\,e^{i\pi/6},
    \end{aligned}
    $

    $
    \begin{aligned}
    \omega &= \dfrac{b}{1-a} \\[1mm]
           &= \dfrac{\dfrac{1 - i\sqrt3}{2}}{\dfrac{1 - i\sqrt3}{4}} \\[1mm]
           &= 2.
    \end{aligned}
    $

    \[
    \boxed{f \text{ est la similitude directe de centre } \Omega(2) \text{, de rapport } \dfrac{\sqrt3}{2} \text{ et d'angle } \dfrac{\pi}{6}.}
    \]

    \item $z' = iz + 1$.

    $
    \begin{aligned}
    a &= i \\[1mm]
      &= e^{i\pi/2}, \qquad |a| = 1,
    \end{aligned}
    $

    $
    \begin{aligned}
    \omega &= \dfrac{b}{1-a} \\[1mm]
           &= \dfrac{1}{1-i} \\[1mm]
           &= \dfrac{1+i}{2}.
    \end{aligned}
    $

    \[
    \boxed{f \text{ est la rotation de centre } \Omega\left(\dfrac{1+i}{2}\right) \text{ et d'angle } \dfrac{\pi}{2}.}
    \]

    \item $z' = z + 3 - i$.

    $
    \begin{aligned}
    a &= 1,
    \end{aligned}
    $

    donc $f$ n'a pas de point invariant.

    \[
    \boxed{f \text{ est la translation de vecteur } \vec{u} \text{ d'affixe } 3 - i.}
    \]

\end{enumerate}

\fbox{\textbf{\exo}} 

Soit $s$ la similitude directe d'écriture complexe $z' = az + b$ avec $a \in \mathbb{C}^*$ et $b \in \mathbb{C}$.
Si $s(A) = A'$ et $s(B) = B'$, alors
\[
\left\{
\begin{aligned}
z_{A'} &= a\,z_A + b \\
z_{B'} &= a\,z_B + b
\end{aligned}
\right.
\]
En soustrayant ces deux égalités, on obtient $z_{A'} - z_{B'} = a\,(z_A - z_B)$, d'où
\[
a = \dfrac{z_{A'} - z_{B'}}{z_A - z_B}
\quad \text{et} \quad
b = z_{A'} - a\,z_A .
\]

\begin{multicols}{2}
\setlength{\columnseprule}{0.1mm} % La largeur de la ligne verticale entre les colonnes
\begin{enumerate}

    \item $A(3+2i)$, $A'(3)$, $B(1)$ et $B'(i)$.

    $
    \begin{aligned}
    a &= \dfrac{z_{A'} - z_{B'}}{z_A - z_B} \\[1mm]
      &= \dfrac{3 - i}{(3+2i) - 1} \\[1mm]
      &= \dfrac{3 - i}{2 + 2i} \\[1mm]
      &= \dfrac{(3 - i)(2 - 2i)}{(2+2i)(2-2i)} \\[1mm]
      &= \dfrac{6 - 6i - 2i + 2i^2}{8} \\[1mm]
      &= \dfrac{4 - 8i}{8} \\[1mm]
      &= \dfrac12 - i,
    \end{aligned}
    $

    $
    \begin{aligned}
    b &= z_{A'} - a\,z_A \\[1mm]
      &= 3 - \left(\dfrac12 - i\right)(3 + 2i) \\[1mm]
      &= 3 - \left(\dfrac32 + i - 3i - 2i^2\right) \\[1mm]
      &= 3 - \left(\dfrac72 - 2i\right) \\[1mm]
      &= -\dfrac12 + 2i.
    \end{aligned}
    $

    \[
    \boxed{z' = \left(\dfrac12 - i\right)z - \dfrac12 + 2i}
    \]

    \item $A(2+i)$, $A'(3+2i)$, $B(2)$ et $B'(3i)$.

    $
    \begin{aligned}
    a &= \dfrac{z_{A'} - z_{B'}}{z_A - z_B} \\[1mm]
      &= \dfrac{(3+2i) - 3i}{(2+i) - 2} \\[1mm]
      &= \dfrac{3 - i}{i} \\[1mm]
      &= \dfrac{(3 - i)(-i)}{i\times(-i)} \\[1mm]
      &= -3i + i^2 \\[1mm]
      &= -1 - 3i,
    \end{aligned}
    $

    $
    \begin{aligned}
    b &= z_{B'} - a\,z_B \\[1mm]
      &= 3i - (-1 - 3i)\times 2 \\[1mm]
      &= 3i + 2 + 6i \\[1mm]
      &= 2 + 9i.
    \end{aligned}
    $

    \[
    \boxed{z' = (-1 - 3i)\,z + 2 + 9i}
    \]

\end{enumerate}
\end{multicols}

\end{document}
